%%% BEGIN fig:dis-twist-graphical
\begin{figure}[htbp]
\centering
% frame of one panel: hadron blob, hard blob, photons, cut (origin at the centre of the hadron blob)
\newcommand{\fbTWframe}{%
  \draw[photon] (-1.65,3.05) -- (-0.95,2.45);
  \draw[photon] (0.95,2.45) -- (1.65,3.05);
  \node[hard blob,ellipse,minimum width=2.1cm,minimum height=0.62cm,inner sep=0pt] at (0,2.35) {};
  \node[font=\small] at (-0.55,2.35) {$H$};
  \draw[black!70,line width=1.5pt] (-2.0,0) -- (2.0,0);
  \node[hadron blob,ellipse,minimum width=2.6cm,minimum height=0.78cm,inner sep=0pt] at (0,0) {};
  \node[font=\small] at (-0.75,0) {$P$};
  \draw[cut] (0,-0.6) -- (0,3.1);}
\begin{tikzpicture}[fd,
  fb lab/.style={font=\footnotesize,align=center,anchor=north}]
  % ======================= (a) twist 2: handbag =======================
  \begin{scope}
    \draw[quark] (-0.55,0.37) -- (-0.55,2.05);
    \draw[quark] (0.55,2.05) -- (0.55,0.37);
    \fbTWframe
    \node[fb lab] at (0,-0.75) {(a)\ 扭度 2（手提包）\\相对权重 $\lambda^0$};
  \end{scope}
  % ======================= (b) one extra transverse gluon =======================
  \begin{scope}[xshift=4.3cm]
    \draw[quark] (-0.85,0.35) -- (-0.85,2.1);
    \draw[gluon] (-0.25,0.39) -- (-0.25,2.04);
    \draw[quark] (0.55,2.05) -- (0.55,0.37);
    \fbTWframe
    \node[font=\scriptsize,anchor=east,inner sep=1pt] at (-0.36,1.66) {$A_T$};
    \node[fb lab] at (0,-0.75) {(b)\ 多一条横向胶子\\$\times\lambda$（扭度 3）};
  \end{scope}
  % ======================= (c) one extra quark pair =======================
  \begin{scope}[xshift=8.6cm]
    \draw[quark] (-0.8,0.35) -- (-0.8,2.1);
    \draw[quark] (-0.35,0.39) -- (-0.35,2.04);
    \draw[quark] (0.35,2.04) -- (0.35,0.39);
    \draw[quark] (0.8,2.1) -- (0.8,0.35);
    \fbTWframe
    \node[fb lab] at (0,-0.75) {(c)\ 多一对夸克\\$\times\lambda^2$（扭度 4）};
  \end{scope}
  % note on scalar-polarized gluons
  \node[anchor=west,align=left,font=\footnotesize,text=blue!50!black] at (10.75,1.35) {纵向极化（K）胶子：\\$\times\lambda^0$，归入\\Wilson 线};
\end{tikzpicture}
\caption{扭度与额外部分子（DIS 强子张量的割图示意，$\lambda\sim\Lambda/Q$）：与硬部分 $H$ 相连的每一条额外的物理极化部分子都带来一个压低因子——每多一条横向胶子压低 $\lambda$，每多一对夸克压低 $\lambda^2$；
纵向极化（K）胶子不受压低，经 Ward 恒等式归入 Wilson 线。非极化单举 DIS 中奇扭度项在无质量微扰论中因手征而消失（\cref{prop:dis-chiral-suppression}），领头的幂次修正来自扭度 4（\cref{prop:dis-twist4}，参见\cref{fig:higher-twist}）。}
\label{fig:dis-twist-graphical}
\end{figure}
%%% END fig:dis-twist-graphical
